\subsection{Gorenstein 环}

定义 1.--- 称环 A 为 Gorenstein 环，如果它是 Noether 的，并且 \(\mathrm{A}_{\mathfrak{m}}\) -模 \(\mathrm{A}_{\mathfrak{m}}\) 对 A 的每个极大理想 \(\mathfrak{m}\) 都具有有限内射维数。

为使 Noether 局部环 A 成为 Gorenstein 环，必须且只需 \(\mathrm{di}_{\mathrm{A}}(\mathrm{A})\) 有限；为使 Noether 环 A 成为 Gorenstein 环，必须且只需对 A 的每个极大理想 m，\(\Lambda_{m}\) 都如此。

命题 10.- 设 A 为 Gorenstein 环；则 A 是 Macaulay 环，且满足 \(\mathrm{di}_{\Lambda}(\mathrm{A}) = \dim(\mathrm{A})\) 。

对 A 的每个极大理想 m，有

\[
\begin{array}{l l} \dim (A _ {m}) & \leqslant \mathrm{di} _ {A _ {m}} (A _ {m}) \\ \mathrm{di} _ {A _ {m}} (A _ {m}) & = \operatorname{prof} (A _ {m}) \\ \operatorname{prof} (A _ {m}) & \leqslant \dim (A _ {m}) \end{array} \quad \begin{array}{l l} & (n ^ {\circ} 6, \text { prop. } 8) \\ & (n ^ {\circ} 6, \text { prop. } 9) \\ & (\S 1, n ^ {\circ} 4, \text { cor. } 2 \text { du   th. } 2); \end{array}
\]

由此推出 A 是 Macaulay 环，并且通过取上确界得到 \(\mathrm{di}_{\mathrm{A}}(\mathrm{A}) = \dim(\mathrm{A})\) (第 2 节, 命题 3)。

于是，使 \(\mathrm{di}_{\mathrm{A}}(\mathrm{A})\) 有限的那些 Noether 环 A 正是维数有限的 Gorenstein 环 (第 2 节, 命题 3)，而使 A-模 A 为内射模的那些 Noether 环正是 Artin Gorenstein 环。

例. 1) 对 Gorenstein 环 A 的每个乘性子集 S，分式环 \(S^{-1}A\) 都是 Gorenstein 环：事实上，设 q 为 \(S^{-1}A\) 的一个极大理想；它形如 \(S^{-1}p\)，其中 p 是 A 的一个与 S 不相交的素理想。设 m 为 A 的包含 p 的一个极大理想；则环 \(B = (S^{-1}A)_q\) 同构于 \(A_p\) (II, § 2, 第 5 节, 命题 11)，从而同构于 \(A_m\) 的一个分式环，因此满足 \(\text{di}_B(B) < +\infty\) (第 2 节, 命题 3 的推论 1)。

\begin{enumerate}
\def\labelenumi{\arabic{enumi})}
\setcounter{enumi}{1}
\item
  设 A 为 Gorenstein 环，J 为 A 的一个由 A-正则序列 x 生成的理想。商环 A/J 是 Gorenstein 环：事实上，对 A 的每个包含 J 的极大理想 m，序列 x 在 \(A_{m}\) 中的像是 \(A_{m}\) -正则的且生成理想 \(J_{m}\)，于是由第 4 节, 命题 7 的推论，\(A_{m}/J_{m}\) 是 Gorenstein 环。
\item
  设 A 为 Noether 局部环，J 为 A 的一个理想，它由 \(\mathfrak{m}_{\mathrm{A}}\) 中元素组成的 A-正则序列生成。若 A/J 是 Gorenstein 环，则 A 亦然 (第 4 节, 命题 7 的推论)。
\item
  设 A 为 Noether 正则局部环；则 A 是 Gorenstein 环。事实上，设 x 为 A 的一个坐标系 (VIII, § 5, 第 1 节, 定义 1)。序列 x 是 A-正则的（同上引文, 第 2 节, 定理 1）且生成理想 m\_A；因此可以应用例 3。
\item
  主理想环的每个商环都是 Gorenstein 环 (例 2)。特别地，域上的每个单生成代数都是 Gorenstein 环。
\end{enumerate}

引理 1.- 设 A 为 Artin 局部环。下列条件等价：

\begin{enumerate}
\def\labelenumi{(\roman{enumi})}
\item
  A 是 Gorenstein 环；
\item
  A-模 A 是内射的；
\item
  \(\kappa_{\mathrm{A}}\) -向量空间 \(\operatorname{Hom}_{\mathrm{A}}(\kappa_{\mathrm{A}}, \mathrm{A})\) 的维数为 1。
\end{enumerate}

回忆 (A, VIII, 第 3.5 页)：A 具有非零的极小理想；这样的理想是单模，从而同构于 \(\kappa_{\mathrm{A}}\) 。于是 \(\kappa_{\mathrm{A}}\) -向量空间 \(\mathrm{Hom}_{\mathrm{A}}(\kappa_{\mathrm{A}}, \Lambda)\) 非零；说它维数为 1，就是说 A 只含一个非零极小理想，它便是 A 的底座 (A, VIII, § 4, 第 6 节)。

(i) 与 (ii) 的等价性已在命题 10 之后证明。设 A-模 A 是内射的。设 \(x\) 与 \(y\) 为 A 的两个被 \(\mathfrak{m}_{\Lambda}\) 零化的非零元素。存在唯一的 A-线性映射 \(\varphi : Ax \to A\) 使 \(\varphi(x) = y\) ；由于 A 是内射的，它可扩张为 \(\Lambda\) 的一个自同态，这蕴含 \(y\) 属于 \(Ax\)，于是得到 (iii)。

反之，设 \(\mathrm{Hom}_{\mathrm{A}}(\kappa_{\mathrm{A}},\Lambda)\) 的维数为 1。设 M 为有限生成 A-模；它具有有限长度，从而具有一个合成列，其诸商都同构于 \(\kappa_{\mathrm{A}}\) 。由此对 M 的长度用归纳法，得到不等式 \(\mathrm{long}_{\mathrm{A}}(\mathrm{Hom}_{\mathrm{A}}(\mathrm{M},\Lambda)) \leqslant \mathrm{long}_{\mathrm{A}}(\mathrm{M})\) 。在扩张模的正合列

\[
0 \to \operatorname{Hom} _ {\mathrm{A}} (\kappa_ {\mathrm{A}}, \mathrm{A}) \longrightarrow \operatorname{Hom} _ {\mathrm{A}} (\mathrm{A}, \mathrm{A}) \stackrel {{\alpha}} {{\longrightarrow}} \operatorname{Hom} _ {\mathrm{A}} (\mathfrak {m} _ {\mathrm{A}}, \mathrm{A}) \longrightarrow \operatorname{Ext} _ {\mathrm{A}} ^ {1} (\kappa_ {\Lambda}, \mathrm{A}) \to 0,
\]

中因此有 \(\operatorname{long}_{\mathrm{A}}(\operatorname{Hom}_{\mathrm{A}}(\mathfrak{m}_{\mathrm{A}},\mathrm{A})) \leqslant \operatorname{long}_{\mathrm{A}}(\mathfrak{m}_{\mathrm{A}}) = \operatorname{long}(\mathrm{A}) - 1 = \operatorname{long}_{\mathrm{A}}(\operatorname{Im}\alpha)\) 。于是 \(\alpha\) 是满射，\(\operatorname{Ext}_{\mathrm{A}}^{1}(\kappa_{\mathrm{A}},\mathrm{A})\) 为零，且 A-模 A 是内射的 (第 3 节, 命题 4)。

引理 2.--- 设 \(\Lambda\) 为满足 \(\mathrm{di}_{\mathbf{A}}(\mathbf{A}) = +\infty\) 的 Noether 局部环。有 \(\operatorname{Ext}_{\mathbf{A}}^{i}(\kappa_{\Lambda}, \mathbf{A}) \neq 0\)（对每个整数 \(i \geqslant \dim(\mathbf{A})\)）。

对 \(\dim(A)\) 用归纳法。当 \(\dim(A) = 0\) 时，\(\mathfrak{m}_A\) 是 \(A\) 的唯一素理想，断言由第 3 节的命题 6 得出。于是设 \(\dim(A) > 0\)，并设 \(\mathfrak{p}\) 为 \(A\) 的一个异于 \(\mathfrak{m}_A\) 的素理想；这时有 \(\dim(A_{\mathfrak{p}}) < \dim(A)\)。若 \(\mathrm{di}_{A_{\mathfrak{p}}}(A_{\mathfrak{p}}) = +\infty\)，则归纳假设蕴含 \(\operatorname{Ext}_{A_{\mathfrak{p}}}^{j}(\kappa(\mathfrak{p}), A_{\mathfrak{p}}) \neq 0\)（对每个整数 \(j \geqslant \dim(A_{\mathfrak{p}})\)）；由 § 1, 第 7 节的引理 3，这蕴含 \(\operatorname{Ext}_{A}^{i}(\kappa_{A}, A) \neq 0\)（对每个整数 \(i \geqslant \dim(A_{\mathfrak{p}}) + \dim(A/\mathfrak{p})\)，特别地对 \(i \geqslant \dim(A)\)）(VIII, § 1, 第 3 节, 命题 8, b))。

尚需处理下列情形：\(\mathrm{di}_{\mathrm{A}}(\mathrm{A})\) 无限，但 \(A_{p}\) 的内射维数对 A 的每个异于 \(m_{A}\) 的素理想 p 都有限。对这样的理想，此时由命题 10 有 \(\mathrm{di}_{\mathrm{A}_{\mathfrak{p}}}(\mathrm{A}_{\mathfrak{p}})=\dim(\mathrm{A}_{\mathfrak{p}})<\dim(\mathrm{A})\)，从而 \(\operatorname{Ext}_{\mathrm{A}_{\mathfrak{p}}}^{i}(\kappa(\mathfrak{p}),\mathrm{A}_{\mathfrak{p}})=0\)（对 \(i\geqslant\dim(\Lambda)\)）。既然 \(\mathrm{di}_{\mathrm{A}}(\mathrm{A})\) 无限，第 3 节的命题 6 便迫使 \(\operatorname{Ext}_{\mathrm{A}}^{i}(\kappa_{\mathrm{A}},\mathrm{A})\neq0\)（对每个整数 \(i\geqslant\dim(\mathrm{A})\)）。

定理 2 (Bass). 设 A 为 Noether 局部环；令 \(d = \dim(A)\)。设 \(\mathbf{x} = \cdot(x_1, \ldots, x_d)\) 为由 \(\mathfrak{m}_A\) 中元素组成的一个极大割序列，并用 \(x\) 表示它生成的理想。下列条件等价：

\begin{enumerate}
\def\labelenumi{(\roman{enumi})}
\item
  A 是 Gorenstein 环；
\item
  有 \(\mathrm{di}_{\mathrm{A}}(\mathrm{A}) = d\)；
\item
  存在整数 \(i>d\) 使 \(\operatorname{Ext}_{\mathrm{A}}^{i}(\kappa_{\mathrm{A}},\mathrm{A})=0\)；
\item
  \(\operatorname{Ext}_{\mathrm{A}}^{i}(\kappa_{\mathrm{A}},\mathrm{A}) = 0\)（对 \(i < d\)），且 \(\kappa_{\mathrm{A}}\) -向量空间 \(\operatorname{Ext}_{\mathrm{A}}^{d}(\kappa_{\mathrm{A}},\mathrm{A})\) 的维数为 1；
\item
  环 A 是 Macaulay 的，且 \(\kappa_{\mathrm{A}}\) -向量空间 \(\operatorname{Hom}_{\Lambda}(\kappa_{\mathrm{A}}, \mathrm{A}/x)\) 的维数为 1；
\item
  序列 \(\mathbf{x}\) 是 A-正则的，且 \(\kappa_{\mathrm{A}}\) -向量空间 \(\operatorname{Hom}_{\mathrm{A}}(\kappa_{\mathrm{A}},\mathrm{A} / x)\) 的维数为 1。
\end{enumerate}

(i)、(ii) 与 (iii) 的等价性由引理 2 及命题 10 得出。若 \(\mathrm{Ext}_{\mathrm{A}}^{i}(\kappa_{\mathrm{A}},\mathrm{A})\) 对每个整数 \(i < d\) 为零，则环 A 是 Macaulay 的 (§ 2, 第 3 节, 命题 3)；若环 A 是 Macaulay 的，则序列 x 是 A-正则的（同上引文, 定理 1）；若序列 x 是 A-正则的，则有 \(\mathrm{Ext}_{\mathrm{A}}^{i}(\kappa_{\mathrm{A}},\mathrm{A}) = 0\)（对 \(i < d\)），且 \(\kappa_{\mathrm{A}}\) -向量空间 \(\mathrm{Ext}_{\mathrm{A}}^{d}(\kappa_{\mathrm{A}},\mathrm{A})\) 与 \(\mathrm{Hom}_{\mathrm{A}}(\kappa_{\mathrm{A}},\mathrm{A}/x)\) 同构 (A, X, 第 166 页, 命题 9)。这就证明了条件 (iv)、(v) 与 (vi) 的等价性。

我们来证明 (i) 蕴含 (v)：若 A 是 Gorenstein 环，则它是 Macaulay 环 (命题 10)。于是序列 x 是 A-正则的 (§ 2, 第 3 节, 命题 4)，从而 A/x 是 Artin Gorenstein 环 (例 2)，故 \(\kappa_{\mathrm{A}}\) -向量空间 \(\mathrm{Hom}_{\mathrm{A}}(\kappa_{\mathrm{A}}, \mathrm{A}/x)\) 的维数为 1。

最后，我们来证明 (vi) 蕴含 (i)：在假设 (vi) 下，环 A/x 是 Gorenstein 环 (引理 1)，从而 A 是 Gorenstein 环 (例 3)。

推论.--- 设 Λ 为维数为 d 的 Noether 局部环。Λ-模 \(\operatorname{Ext}_{\mathrm{A}}^{d}(\mathbf{k}_{\mathrm{A}},\Lambda)\) 非零。

若 \(\Lambda\) 是 Gorenstein 环，这由定理 2 得出；否则由引理 2 得出。

命题 11.--- 设 A 为 Noether 环。下列条件等价：

\begin{enumerate}
\def\labelenumi{(\roman{enumi})}
\item
  A 是 Gorenstein 环；
\item
  对 A 的每个素理想 \(\mathfrak{p}\)，\(\kappa(\mathfrak{p})\) -向量空间 \(\operatorname{Ext}_{\Lambda_{\mathfrak{p}}}^{i}(\kappa(\mathfrak{p}),\mathrm{A}_{\mathfrak{p}})\) 在 \(i \neq \operatorname{ht}(\mathfrak{p})\) 时为零，而在 \(i = \operatorname{ht}(\mathfrak{p})\) 时维数为 1。
\item
  对 A 的每个极大理想 \(\mathfrak{m}\)，存在整数 \(i > \mathrm{ht}(\mathfrak{m})\) 使 \(\operatorname{Ext}_{\Lambda_{\mathfrak{m}}}^{i}(\mathrm{A}/\mathfrak{m},\mathrm{A}_{\mathfrak{m}}) = 0\) 。
\end{enumerate}

(i)⇒(ii)：这由把定理 2 应用于 Gorenstein 局部环 A \(_{p}\) 得出 (例 1)。

(ii)⇒(iii)：这是平凡的。

(iii)⇒(i)：在假设 (iii) 下，对 A 的每个极大理想 m，A \(_{m}\) 都是 Gorenstein 环 (定理 2)，故 A 是 Gorenstein 的。

